% Fragment for the article. The accompanying methods and source tables are
% documented in FLOQUET_PILOT_RESULTS.md. The table uses booktabs.
\paragraph{Local stability of periodic RF orbits.}
We found 66 noncentral periodic solutions on six branches of the full
time-dependent N48 field over \(0.40\leq\lambda_4\leq0.72\). Each branch
was sampled at 11 coupling values. For every computed orbit, at least one
Floquet multiplier lay outside the unit circle. Native P4 calculations
on both meshes confirmed the local classification for three selected orbits
(Table~\ref{tab:floquet-p4}).

\begin{table}[ht]
\centering
\caption{Largest Floquet multiplier modulus for three periodic orbits.
The orbit period is expressed in RF periods.}
\label{tab:floquet-p4}
\begin{tabular}{rrrrr}
\toprule
\(\lambda_4\) & Period & N48 & P4 h35 & P4 h50 \\
\midrule
0.45266 & 28 & 1.0125796 & 1.0125859 & 1.0124558 \\
0.56511 & 4  & 7228.3484 & 7228.3485 & 7228.3485 \\
0.67756 & 20 & 1.2582715 & 1.2582714 & 1.2582620 \\
\bottomrule
\end{tabular}
\end{table}

All six nearby and nominal launches in the separate N48 neighborhood
test remained inside the aperture for 1000 RF periods. This demonstrates
that finite-time capture can coexist with local linear instability of
periodic orbits. No sampled branch changed its linear stability class
across the investigated coupling interval. The finite-time escape
mechanism remains unresolved.
